\scanpage{93}

% L1–2
on peut supposer ces conditions satisfaites, \deleted{\unsure{donc}} \inserted{donc aussi (9)} \inserted{donc aussi (10) vrai.} \deleted{On \unsure{aura}} \inserted{Notons} qu'alors \deleted{l'indépendance de} la condition

% L3–6
$s_{n+1}$ \deleted{\unsure{$f$|$\notin$}} libre mod $F_n$ équivaut à la condition \deleted{libre mod $F_n$} (5), \deleted{Maintenant on sait déjà que $a'_n\mid F_{n-1}=0$, donc} \inserted{Mais on sait qu'elle} équivaut à $\langle s_n,a'_n\rangle$ \deleted{$\in$ \unsure{inversible|inv.}}, et $\langle s_n,a'_n\rangle=\lambda_n$, OK.

% L7–9, diagonal cancellation begins
\deleted{Il reste, pour compléter la récurrence, à établir (8) pour $i=n+1$, or on \unsure{aura|a}}

% L10, cancelled display
\[
\deleted{s_{n+1}\overset{\text{déf}}{=}\tau_n s_n=s_n+\langle s_n,a'_n\rangle a_n\equiv\langle s_n,a'_n\rangle a_n\qquad(F_n)}
\]

% L11–14
\deleted{et comme $\underset{\in F_n}{s_n}$ \unsure{$s_n\equiv$} $\lambda_{n-1}a_{n-1}\langle s_{n-1},a'_n\rangle a_n\equiv\langle s_n,a'_n\rangle a_{n-1}$ mod $F_{n-1}\subset F_n$ donc et \unsure{fait|mod} $F_n=\mathcal O_S(s_0,a_0,\ldots,a_{n-1})$ mod $F_n$, \unsure{on aura} \unsure{déjà acquis}}

% L13–16, nested workings
\[
\begin{aligned}
&\deleted{\left\{\begin{aligned}
s_n&\equiv\deleted{\unsure{s_{n-1}+}}\lambda_{n-1}a_{n-1}\qquad(F_{n-1})\\
a'_n&\mid F_{n-1}=0
\end{aligned}\right.}\\
&\deleted{s_{n+1}\equiv\langle s_n,a'}\\
&\deleted{s_{n+1}\equiv\lambda_{n+1}\lambda_{n-1}a_{n-1}}\qquad
\deleted{\overset{\Vert}{\lambda_n a_n}\quad\text{par (9)}}\\
&\deleted{\text{d'où }\langle s_n,a'_n\rangle=\unsure{\lambda_n.}}\qquad\deleted{\text{OK.}}
\end{aligned}
\]

% L15–17, left
% slanted
\margin{L'assertion de stabilité évidente\ldots}

% L18–19
\deleted{\uline{Corollaire} Pour fibré vectoriel ($\mathcal O$ Mod. l.l. de \unsure{r.f.|t.f.}) et conditions fibre par fibre.}

% L20–22
\deleted{\uline{Reformulation} \deleted{\unsure{aussi}} Cas où $E$ loc. libre de rang $n+2$. Conditions équivalentes sur le système $(\overset{\text{section}}{s_0},\overbrace{\unsure{\tau_0|\mathrm T_0},\ldots,\tau_n}^{\text{ps.réflexions}})$, définissant les $s_i$}

% L23–24, bottom upside down
% margin note written sideways (upside down)
\margin{$a'_i\otimes \underline{a_i}$\\ $\underline{a'_{i+1}}\otimes a_{i+1}$}

% L23–24, bottom upside down, right of products
% upside down
\diagram{\begin{tikzcd}[ampersand replacement=\&,row sep=small,column sep=small]
\deleted{\unsure{\xi}} \& a_0 \arrow[r,dash]\arrow[dr,dash,"{a'_1}"] \& a_n\\
\&\&{}
\end{tikzcd}}

% archive sheet number: 92
